% !PRINTED 333
\weberpage{333}
%<<CONT

\noindent 3.\ \gesperrt{Eine ternäre cubische Form $f(x_1,x_2,x_3)$, deren Discriminante von Null verschieden ist, lässt sich auf vier verschiedene Arten in die canonische Form}
\[
\varphi(y_1,y_2,y_3) = h(y_1^3+y_2^3+y_3^3) + 6\,m\,y_1y_2y_3
\]
\gesperrt{transformiren.}

Da auf jeder der Linien $y_1=0$, $y_2=0$, $y_3=0$ drei Inflexionspunkte liegen, so ist das Problem der Inflexionspunkte wesentlich identisch mit dem der Transformation der cubischen Form $f$ auf die canonische Form, mit dem wir uns zunächst beschäftigen wollen $^1$).

\begin{center}
\S.\,90.

\gesperrt{Transformation der cubischen Form auf die canonische Form.}
\end{center}

Um die Transformation der cubischen ternären Form auf die canonische Form durchzuführen, d. h. auf die Lösung gewisser algebraischer Gleichungen zurückzuführen, müssen wir zunächst die Covarianten für die canonische Form bilden.

Es sei also
\[
\varphi(y) = h(y_1^3+y_2^3+y_3^3) + 6\,m\,y_1y_2y_3,
\]
wofür wir zur Abkürzung auch setzen:
\weq{1}{\varphi(y) = h\,\Sigma y_1^3 + 6\,m\,y_1y_2y_3,}
indem wir unter dem Zeichen $\Sigma$ eine Summe über drei durch cyklische Vertauschung der Variablen $y_1,y_2,y_3$ aus dem ersten gebildete Glieder verstehen.

Die \gesperrt{Hesse'sche} Covariante wollen wir zur Vereinfachung der Formeln mit einem geeigneten numerischen Factor multipliciren, und demnach
\weq{2}{\Delta(x) = \frac{1}{6^3}\left|\begin{array}{ccc} f_{11}, & f_{12}, & f_{13} \\ f_{21}, & f_{22}, & f_{23} \\ f_{31}, & f_{32}, & f_{33} \end{array}\right|}

\medskip
\noindent\rule{4cm}{0.4pt}

\smallskip
\noindent $^1$) Aus der Literatur über die Wendepunkte der Curven dritter Ordnung und über die Theorie der ternären cubischen Formen erwähnen wir ausser den älteren Arbeiten von \gesperrt{Newton} und \gesperrt{Mac Laurin}: \gesperrt{Hesse}, \glqq Ueber die Elimination etc.\grqq{} und \glqq Ueber die Wendepunkte der Curven dritter Ordnung\grqq{}, \gesperrt{Crelle's} Journal, Bd. 28 (1844). \gesperrt{Aronhold}, Theorie der homogenen Functionen 3\textsuperscript{ten} Grades von drei Veränderlichen, \gesperrt{Crelle's} Journal, Bd. 55 (1858). \gesperrt{Salmon}, Higher plane curves (deutsch von \gesperrt{Fiedler}).
%CONT>>
